% Source problem: public/notes/物工/問題/2015ap_exam.pdf, PDF page 7.
\documentclass[11pt,a4paper]{article}
\usepackage[margin=24mm]{geometry}
\usepackage{fontspec}
\usepackage{xeCJK}
\setCJKmainfont{HaranoAjiMincho-Regular.otf}
\setCJKsansfont{HaranoAjiGothic-Medium.otf}
\usepackage{amsmath,amssymb}
\setlength{\parindent}{1em}
\setlength{\parskip}{0.35em}
\renewcommand{\figurename}{図}
\title{物理工学専攻 院試解答集\\2015年 物理学II 第1問}
\author{政岡凜太郎}
\date{}
\begin{document}
\maketitle

\section*{[1] 数演算子と消滅演算子}
交換子の積の公式$[AB,C]=A[B,C]+[A,C]B$を用いると
\begin{equation}
  [\hat N,\hat a]=[\hat a^\dagger\hat a,\hat a]
  =\hat a^\dagger[\hat a,\hat a]+[\hat a^\dagger,\hat a]\hat a
  =-\hat a.
\end{equation}
このため
\begin{equation}
  \hat N(\hat a|n+1\rangle)
  =\hat a\hat N|n+1\rangle-\hat a|n+1\rangle
  =n\hat a|n+1\rangle.
\end{equation}
従って$\hat a|n+1\rangle$は固有値$n$の固有状態に比例する。規格化された固有状態の位相を適当に選び、$\hat a|n+1\rangle=c_n|n\rangle$（$c_n\ge0$）と書く。係数はノルムの二乗から
\begin{equation}
  c_n^2=\langle n+1|\hat a^\dagger\hat a|n+1\rangle=n+1
\end{equation}
なので、
\begin{equation}
  \boxed{\hat a|n+1\rangle=\sqrt{n+1}|n\rangle.}
\end{equation}
同様に
\begin{equation}
  \|\hat a|0\rangle\|^2=\langle0|\hat N|0\rangle=0
\end{equation}
だから$\boxed{\hat a|0\rangle=0}$である。

\section*{[2] コヒーレント状態の固有値}
互いに直交規格化された数状態を使うと
\begin{equation}
  \langle\Psi(\alpha)|\Psi(\alpha)\rangle
  =e^{-|\alpha|^2}\sum_{n=0}^{\infty}\frac{|\alpha|^{2n}}{n!}=1.
\end{equation}
消滅演算子を級数に作用させ、[1]を使って添字をずらすと
\begin{align}
  \hat a|\Psi(\alpha)\rangle
  &=e^{-|\alpha|^2/2}\sum_{n=1}^{\infty}
    \frac{\alpha^n}{\sqrt{n!}}\sqrt n\,|n-1\rangle\\
  &=\alpha e^{-|\alpha|^2/2}\sum_{j=0}^{\infty}
    \frac{\alpha^j}{\sqrt{j!}}|j\rangle
   =\boxed{\alpha|\Psi(\alpha)\rangle}.
\end{align}

\section*{[3] エネルギー期待値}
[2]の固有値式と規格化を、$\hat N=\hat a^\dagger\hat a$ の期待値に用いると
\begin{equation}
 \langle\hat N\rangle
 =\|\hat a|\Psi(\alpha)\rangle\|^2
 =|\alpha|^2\langle\Psi(\alpha)|\Psi(\alpha)\rangle
 =|\alpha|^2.
\end{equation}
従って
\begin{equation}
  \boxed{\langle\hat H\rangle
  =\hbar\omega\left(|\alpha|^2+\frac12\right).}
\end{equation}

\section*{[4] 平均値と不確定性}
生成・消滅演算子の定義を足し引きすると
\begin{equation}
  \hat x=\sqrt{\frac{\hbar}{2m\omega}}(\hat a+\hat a^\dagger),
  \qquad
  \hat p=-i\sqrt{\frac{m\hbar\omega}{2}}(\hat a-\hat a^\dagger).
\end{equation}
固有状態の性質から$\langle\hat a\rangle=\alpha$、$\langle\hat a^\dagger\rangle=\alpha^*$なので
\begin{equation}
  \boxed{\langle\hat x\rangle=\sqrt{\frac{2\hbar}{m\omega}}\,\Re\alpha},
  \qquad
  \boxed{\langle\hat p\rangle=\sqrt{2m\hbar\omega}\,\Im\alpha}.
\end{equation}
また$\langle\hat a^2\rangle=\alpha^2$、$\langle\hat a^{\dagger2}\rangle=(\alpha^*)^2$、$\langle\hat a^\dagger\hat a\rangle=|\alpha|^2$、$\langle\hat a\hat a^\dagger\rangle=|\alpha|^2+1$より
\begin{align}
  \langle\hat x^2\rangle
  &=\frac{\hbar}{2m\omega}\left[(\alpha+\alpha^*)^2+1\right],\\
  \langle\hat p^2\rangle
  &=-\frac{m\hbar\omega}{2}\left[(\alpha-\alpha^*)^2-1\right].
\end{align}
平均値の二乗をそれぞれ引くと、
\begin{equation}
  \boxed{\Delta x=\sqrt{\frac{\hbar}{2m\omega}}},
  \qquad
  \boxed{\Delta p=\sqrt{\frac{m\hbar\omega}{2}}},
  \qquad
  \boxed{\Delta x\Delta p=\frac{\hbar}{2}}.
\end{equation}
各標準偏差は正であり、積はハイゼンベルクの下限に等しい。

\section*{[5] 時間発展}
数状態ごとの時間発展は
\begin{equation}
  e^{-i\hat Ht/\hbar}|n\rangle
  =e^{-i\omega t(n+1/2)}|n\rangle.
\end{equation}
これを状態の級数に適用すると
\begin{align}
  e^{-i\hat Ht/\hbar}|\Psi(\alpha_0)\rangle
  &=e^{-i\omega t/2}e^{-|\alpha_0|^2/2}
    \sum_{n=0}^{\infty}\frac{(\alpha_0e^{-i\omega t})^n}{\sqrt{n!}}|n\rangle\\
  &=e^{-i\omega t/2}|\Psi(\alpha_0e^{-i\omega t})\rangle.
\end{align}
従って
\begin{equation}
  \boxed{\alpha(t)=\alpha_0e^{-i\omega t}
  =Ae^{i(\theta-\omega t)}}.
\end{equation}
前の因子$e^{-i\omega t/2}$は全ての数状態に共通する全体位相であり、期待値には影響しない。

\section*{[6] 古典的な運動と準古典状態}
時刻$t$で$\alpha(t)=Ae^{i(\theta-\omega t)}$を位置・運動量の平均式に代入すると
\begin{align}
  \boxed{\langle\hat x\rangle_t
  =\sqrt{\frac{2\hbar}{m\omega}}A\cos(\omega t-\theta)},\\
  \boxed{\langle\hat p\rangle_t
  =-\sqrt{2m\hbar\omega}A\sin(\omega t-\theta)}.
\end{align}
これらは古典調和振動子の位置・運動量の時間変化そのものである。振幅は$A$に比例する一方、$\Delta x$と$\Delta p$は$A$に依存しない。それぞれの振動振幅に対する標準偏差の比は$1/(2A)$であり、$A\gg1$では波束の広がりが古典軌道の振幅より十分小さくなる。この意味で準古典状態と呼ばれる。

\end{document}
